\documentclass[11pt,a4paper]{report}

\usepackage[utf8]{inputenc}
\usepackage[T1]{fontenc}
\usepackage{lmodern}
\usepackage[margin=1in]{geometry}
\usepackage{graphicx}
\usepackage{xcolor}
\usepackage{listings}
\usepackage{hyperref}
\usepackage{booktabs}
\usepackage{enumitem}
\usepackage{fancyhdr}
\usepackage{titlesec}
\usepackage{tocloft}
\usepackage{skak}           % Chess diagrams and notation
\usepackage{xskak}          % Extended chess support
\usepackage{chessboard}     % Chess board diagrams
\usepackage{amsmath}
\usepackage{algorithm}
\usepackage{algpseudocode}

% Colors
\definecolor{codebackground}{RGB}{245,245,245}
\definecolor{codecomment}{RGB}{0,128,0}
\definecolor{codekeyword}{RGB}{0,0,180}
\definecolor{codestring}{RGB}{163,21,21}
\definecolor{linkcolor}{RGB}{0,80,160}
\definecolor{chaptercolor}{RGB}{40,60,100}

% Hyperref setup
\hypersetup{
    colorlinks=true,
    linkcolor=linkcolor,
    urlcolor=linkcolor,
    citecolor=linkcolor,
    pdftitle={Chess Engine Documentation},
    pdfauthor={Your Name},
}

% Code listing style
\lstdefinestyle{codestyle}{
    backgroundcolor=\color{codebackground},
    basicstyle=\ttfamily\small,
    breakatwhitespace=false,
    breaklines=true,
    captionpos=b,
    commentstyle=\color{codecomment}\itshape,
    keywordstyle=\color{codekeyword}\bfseries,
    stringstyle=\color{codestring},
    numberstyle=\tiny\color{gray},
    numbers=left,
    numbersep=8pt,
    showspaces=false,
    showstringspaces=false,
    showtabs=false,
    tabsize=4,
    frame=single,
    framerule=0.5pt,
    rulecolor=\color{gray!50},
}
\lstset{style=codestyle}

% Chapter and section styling
\titleformat{\chapter}[display]
    {\normalfont\huge\bfseries\color{chaptercolor}}
    {\chaptertitlename\ \thechapter}{20pt}{\Huge}
\titlespacing*{\chapter}{0pt}{0pt}{30pt}

% Header and footer
\pagestyle{fancy}
\fancyhf{}
\fancyhead[L]{\leftmark}
\fancyhead[R]{\thepage}
\renewcommand{\headrulewidth}{0.4pt}

% Chess move formatting
\newcommand{\move}[1]{\texttt{#1}}

% Function/method documentation
\newcommand{\function}[4]{%
    \paragraph{\texttt{#1}} \hfill \textit{#2}\\
    \textbf{Parameters:} #3\\
    \textbf{Returns:} #4
}

% API endpoint
\newcommand{\apiendpoint}[3]{%
    \noindent\fbox{\texttt{#1}} \texttt{#2}\\
    #3\vspace{0.5em}
}

% Note box
\newcommand{\notebox}[1]{%
    \begin{center}
    \fcolorbox{gray}{gray!10}{%
        \parbox{0.9\textwidth}{\textbf{Note:} #1}
    }
    \end{center}
}

% Warning box
\newcommand{\warningbox}[1]{%
    \begin{center}
    \fcolorbox{red!70}{red!10}{%
        \parbox{0.9\textwidth}{\textbf{Warning:} #1}
    }
    \end{center}
}

%------------------------------------------------------------------------------
% Document Information
%------------------------------------------------------------------------------
\title{%
    \vspace{-1cm}
    \rule{\textwidth}{1pt}\\[0.5cm]
    {\Huge\bfseries First Move Robotics}\\[0.3cm]
    {\Large Software Documentation}\\[0.3cm]
    \rule{\textwidth}{1pt}\\[1cm]
    {\large Version 1.0.0}
}
\author{Strydr Silverberg\\strydr\_silverberg@mines.edu}
\date{\today}

\begin{document}

\maketitle
\thispagestyle{empty}

\vfill
\begin{center}
    \textit{A high-performance chess engine written in C++}
\end{center}
\vfill

\newpage
\tableofcontents
\newpage

%==============================================================================
\chapter{Introduction}
%==============================================================================

\section{Overview}

This document provides comprehensive documentation for the \textbf{Kirin Chess Engine}. Kirin is a chess engine written in C++ and an integral part of the \textbf{First Move Robotics} chess board project.

\section{Features}

\begin{itemize}
    \item Feature 1: Brief description
    \item Feature 2: Brief description
    \item Feature 3: Brief description
\end{itemize}

\section{System Requirements}

\begin{table}[h]
\centering
\begin{tabular}{ll}
\toprule
\textbf{Component} & \textbf{Requirement} \\
\midrule
Operating System & Windows / Linux / macOS \\
Memory & Minimum 256 MB RAM \\
Processor & Any modern x86-64 CPU \\
Dependencies & List dependencies here \\
\bottomrule
\end{tabular}
\caption{Minimum system requirements}
\end{table}

%==============================================================================
\chapter{Installation}
%==============================================================================

\section{Building from Source}

Clone the repository and build:

\begin{lstlisting}[language=bash]
git clone https://github.com/username/chess-engine.git
cd chess-engine
make
\end{lstlisting}

\section{Pre-built Binaries}

Download the latest release from the releases page.

\section{Configuration}

The engine can be configured via the \texttt{config.ini} file:

\begin{lstlisting}
[Engine]
hash_size = 256
threads = 4
ponder = true
\end{lstlisting}

%==============================================================================
\chapter{Architecture}
%==============================================================================

\section{Overview}

Describe the high-level architecture of your engine here.

\section{Board Representation}

\subsection{Data Structures}

Describe how the board is represented (bitboards, mailbox, 0x88, etc.).

\begin{lstlisting}[language=C++]
// Example bitboard representation
struct Position {
    uint64_t pieces[2][6];  // [color][piece_type]
    uint64_t occupied[2];   // [color]
    uint64_t all_pieces;
    int side_to_move;
    int castling_rights;
    int en_passant_square;
    int halfmove_clock;
    int fullmove_number;
};
\end{lstlisting}

\subsection{Chess Position Example}

Here is an example starting position:

\begin{center}
\chessboard[setfen=rnbqkbnr/pppppppp/8/8/8/8/PPPPPPPP/RNBQKBNR w KQkq - 0 1]
\end{center}

After \move{1.e4 e5 2.Nf3}:

\begin{center}
\chessboard[setfen=rnbqkbnr/pppp1ppp/8/4p3/4P3/5N2/PPPP1PPP/RNBQKB1R b KQkq - 1 2]
\end{center}

\section{Move Generation}

Describe your move generation approach.

\subsection{Legal Move Generation}

\begin{algorithm}
\caption{Legal Move Generation}
\begin{algorithmic}[1]
\Procedure{GenerateMoves}{position}
    \State moves $\gets$ empty list
    \For{each piece of side to move}
        \State Generate pseudo-legal moves for piece
        \For{each pseudo-legal move}
            \If{move does not leave king in check}
                \State Add move to moves
            \EndIf
        \EndFor
    \EndFor
    \State \Return moves
\EndProcedure
\end{algorithmic}
\end{algorithm}

%==============================================================================
\chapter{Search Algorithm}
%==============================================================================

\section{Overview}

Describe your search algorithm (Alpha-Beta, MCTS, etc.).

\section{Alpha-Beta Search}

\begin{algorithm}
\caption{Alpha-Beta with Negamax}
\begin{algorithmic}[1]
\Function{AlphaBeta}{position, depth, $\alpha$, $\beta$}
    \If{depth = 0}
        \State \Return \Call{Evaluate}{position}
    \EndIf
    \For{each move in \Call{GenerateMoves}{position}}
        \State \Call{MakeMove}{position, move}
        \State score $\gets$ $-$\Call{AlphaBeta}{position, depth$-$1, $-\beta$, $-\alpha$}
        \State \Call{UnmakeMove}{position, move}
        \If{score $\geq \beta$}
            \State \Return $\beta$ \Comment{Beta cutoff}
        \EndIf
        \If{score $> \alpha$}
            \State $\alpha \gets$ score
        \EndIf
    \EndFor
    \State \Return $\alpha$
\EndFunction
\end{algorithmic}
\end{algorithm}

\section{Search Enhancements}

\subsection{Transposition Table}

Describe your transposition table implementation.

\subsection{Move Ordering}

\begin{enumerate}
    \item Hash move (from transposition table)
    \item Winning captures (MVV-LVA)
    \item Killer moves
    \item History heuristic
    \item Quiet moves
\end{enumerate}

\subsection{Pruning Techniques}

\begin{itemize}
    \item Null Move Pruning
    \item Late Move Reductions (LMR)
    \item Futility Pruning
    \item Delta Pruning (in quiescence)
\end{itemize}

%==============================================================================
\chapter{Evaluation}
%==============================================================================

\section{Overview}

Describe your evaluation function approach.

\section{Material}

\begin{table}[h]
\centering
\begin{tabular}{lc}
\toprule
\textbf{Piece} & \textbf{Value (centipawns)} \\
\midrule
Pawn & 100 \\
Knight & 320 \\
Bishop & 330 \\
Rook & 500 \\
Queen & 900 \\
\bottomrule
\end{tabular}
\caption{Piece values}
\end{table}

\section{Positional Factors}

\subsection{Piece-Square Tables}

\begin{lstlisting}[language=C++]
// Example: Pawn piece-square table (middlegame)
const int pawn_psqt[64] = {
     0,  0,  0,  0,  0,  0,  0,  0,
    50, 50, 50, 50, 50, 50, 50, 50,
    10, 10, 20, 30, 30, 20, 10, 10,
     5,  5, 10, 25, 25, 10,  5,  5,
     0,  0,  0, 20, 20,  0,  0,  0,
     5, -5,-10,  0,  0,-10, -5,  5,
     5, 10, 10,-20,-20, 10, 10,  5,
     0,  0,  0,  0,  0,  0,  0,  0
};
\end{lstlisting}

\section{King Safety}

Describe king safety evaluation.

\section{Pawn Structure}

Describe pawn structure evaluation: doubled pawns, isolated pawns, passed pawns, etc.

%==============================================================================
\chapter{UCI Protocol}
%==============================================================================

\section{Overview}

The engine communicates via the Universal Chess Interface (UCI) protocol.

\section{Supported Commands}

\subsection{Standard UCI Commands}

\apiendpoint{IN}{uci}{Initialize UCI mode and request engine identification.}

\apiendpoint{IN}{isready}{Check if engine is ready.}

\apiendpoint{IN}{position [fen | startpos] moves ...}{Set up a position.}

\apiendpoint{IN}{go [options]}{Start calculating. Options include:}
\begin{itemize}[noitemsep]
    \item \texttt{depth n} -- search to depth n
    \item \texttt{movetime ms} -- search for ms milliseconds
    \item \texttt{wtime/btime ms} -- time remaining
    \item \texttt{infinite} -- search until stopped
\end{itemize}

\apiendpoint{IN}{stop}{Stop calculating.}

\apiendpoint{IN}{quit}{Exit the engine.}

\subsection{Engine Options}

\begin{lstlisting}
setoption name Hash value 256
setoption name Threads value 4
setoption name Ponder value true
\end{lstlisting}

%==============================================================================
\chapter{API Reference}
%==============================================================================

\section{Core Classes}

\subsection{Position}

\function{Position()}{Constructor}{None}{New Position object}

\function{set\_fen(fen: string)}{Method}{fen -- FEN string}{bool -- success}

\function{make\_move(move: Move)}{Method}{move -- Move to make}{bool -- success}

\subsection{Search}

\function{search(position: Position, limits: SearchLimits)}{Function}{position, limits}{SearchResult}

\section{Utility Functions}

\function{perft(position: Position, depth: int)}{Function}{position, depth}{uint64 -- node count}

%==============================================================================
\chapter{Testing}
%==============================================================================

\section{Perft Testing}

Perft (performance test) validates move generation:

\begin{lstlisting}[language=bash]
./engine perft 6
\end{lstlisting}

\begin{table}[h]
\centering
\begin{tabular}{cr}
\toprule
\textbf{Depth} & \textbf{Nodes} \\
\midrule
1 & 20 \\
2 & 400 \\
3 & 8,902 \\
4 & 197,281 \\
5 & 4,865,609 \\
6 & 119,060,324 \\
\bottomrule
\end{tabular}
\caption{Perft results from starting position}
\end{table}

\section{Test Positions}

Include strategic and tactical test suites (WAC, ECM, STS, etc.).

%==============================================================================
\chapter{Performance}
%==============================================================================

\section{Benchmarks}

Document performance metrics here.

\section{Optimization Techniques}

\begin{itemize}
    \item Magic bitboards for sliding piece attacks
    \item SIMD instructions where applicable
    \item Cache-friendly data structures
    \item Efficient hash table implementation
\end{itemize}

%==============================================================================
\chapter{Contributing}
%==============================================================================

\section{Development Setup}

Instructions for setting up a development environment.

\section{Code Style}

Document your code style guidelines.

\section{Pull Request Process}

\begin{enumerate}
    \item Fork the repository
    \item Create a feature branch
    \item Make changes and add tests
    \item Submit a pull request
\end{enumerate}

%==============================================================================
\appendix
\chapter{Glossary}
%==============================================================================

\begin{description}
    \item[Alpha-Beta] A search algorithm that prunes branches that cannot affect the final result.
    \item[Bitboard] A 64-bit integer representing piece positions on the board.
    \item[Centipawn] One hundredth of a pawn; standard unit for evaluation scores.
    \item[FEN] Forsyth-Edwards Notation; a standard for describing chess positions.
    \item[Perft] Performance test; counts all leaf nodes at a given depth.
    \item[Ply] A half-move; one move by either white or black.
    \item[UCI] Universal Chess Interface; standard protocol for chess engines.
\end{description}

%==============================================================================
\chapter{References}
%==============================================================================

\begin{itemize}
    \item Chess Programming Wiki: \url{https://www.chessprogramming.org/}
    \item UCI Protocol Specification
    \item Stockfish source code: \url{https://github.com/official-stockfish/Stockfish}
\end{itemize}

%%%%%%%%%%%%%%%%%%%%%%%%%%%%%%%%%%%%%%%%%%%%%%%%%%%%%%%%%%%%%%%%%%%%%%%%%%%%%%%
\end{document}
%%%%%%%%%%%%%%%%%%%%%%%%%%%%%%%%%%%%%%%%%%%%%%%%%%%%%%%%%%%%%%%%%%%%%%%%%%%%%%%
